\section{Experimental details and reproduction}

\subsection{Code and data availability}
\label{sec:availability}

Every number in this paper regenerates from the code released with it (the supplementary
material during review). The release holds the autograd engine, the transformer, the analysis
toolkit, the kernels that produced each sweep, and \texttt{scripts/reproduce.sh}. That script
re-runs every analysis from the saved artifacts, with pinned dependencies and fixed seeds. It
does not train a model (Section~\ref{sec:repro-procedure}). Saved artifacts carry a provenance stamp
recording the git SHA, the dtype, the account whose compute produced them and the exact
configuration. \texttt{test\_paper\_numbers.py} re-derives every load-bearing number in this
manuscript, captions included, from those artifacts rather than from the prose, and it runs as
part of the self-check suite. Two limitations of that guarantee are stated in
Section~\ref{sec:limits}. $68$ artifacts from the three earliest sweeps carry no usable SHA,
and their code version is recovered by argument rather than by stamp. Figures are regenerated
rather than stored, so a figure is only as reproducible as the artifact beneath it.

\subsection{The sweeps}
\label{sec:sweeps}

% STE descriptive
Tables~\ref{tab:sweeps} and~\ref{tab:sweep-files} list every sweep that supplies a result. The
first table gives the configuration of each sweep. The second table gives the script that
trained it, the directory of its artifacts, the script that analyses it and its compute.
\texttt{scripts/sweep\_table.py} writes both tables from the artifacts and the run logs.
\texttt{test\_paper\_numbers.py} checks that the tables here are equal to its output.

\begin{table}[ht]
  \centering
  \footnotesize
  \caption{\textbf{The configuration of each sweep.} ``Engine'' is our from-scratch
    implementation and ``PyTorch'' is the reference implementation. Every sweep uses AdamW
    with learning rate $10^{-3}$ and weight decay $1.0$, full-batch. A moduli cell with more
    than four moduli gives their count and range. Table~\ref{tab:moduli} lists each modulus.}
  \label{tab:sweeps}
  \begin{tabular}{llllllll}
    \toprule
    sweep & code & precision & moduli & seeds & steps & train fraction & platform \\
    \midrule
    k01 & PyTorch & float32 & $6$, $113$ to $165$ & $0$ & $25{,}000$ & $0.3$ & Kaggle T4 \\
    k02 & PyTorch & float32 & $6$, $113$ to $165$ & $0$--$4$ & $40{,}000$ & $0.3$ & Kaggle T4 \\
    k03 & PyTorch & float32 & $6$, $113$ to $165$ & $0$--$4$ & $40{,}000$ & $0.3$ & Kaggle T4 \\
    k04 & PyTorch & float32 & $12$, $49$ to $169$ & $0$--$2$ & $40{,}000$ & $0.3$ & Kaggle T4 \\
    k05 & PyTorch & float32 & $49$, $54$, $63$ & $0$--$2$ & $40{,}000$ & $0.3$, $0.5$, $0.8$ & Kaggle T4 \\
    k06 & PyTorch & float32 & $119$, $120$, $125$, $165$ & $0$--$2$ & $120{,}000$ & $0.3$ & Kaggle T4 \\
    k07 & PyTorch & float32 & $113$ & $0$--$4$ & $40{,}000$ & $0.3$ & Kaggle T4 \\
    k08 & PyTorch & float32 & $113$ & $0$--$7$ & $40{,}000$ & $0.3$ & Kaggle T4 \\
    k09 & PyTorch & float32 & $5$, $91$ to $131$ & $0$--$2$ & $40{,}000$ & $0.3$ & Kaggle T4 \\
    N7 & engine & float64 & $113$, $119$, $121$, $125$ & $0$--$2$ & $40{,}000$ & $0.3$ & CPU \\
    N9 & engine & float32 & $113$ & $0$--$2$ & $40{,}000$ & $0.3$ & CPU \\
    O18 & PyTorch & float32 & $113$ & $0$--$2$ & $40{,}000$ & $0.3$ & CPU \\
    O18, f64 & PyTorch & float64 & $113$ & $0$--$2$ & $40{,}000$ & $0.3$ & CPU \\
    O18-T & engine & float64 & $113$ & $0$ & $40{,}000$ & $0.3$ & CPU \\
    I1 & engine & float64 & $113$ & $0$--$2$ & $40{,}000$ & $0.3$ & CPU \\
    I1b & engine & float64 & $121$ & $0$--$2$ & $40{,}000$ & $0.3$ & CPU \\
    Gate 1 & engine & float64 & $113$ & $0$ & $25{,}000$ & $0.3$ & CPU \\
    \bottomrule
  \end{tabular}
\end{table}

\begin{table}[ht]
  \centering
  \footnotesize
  \caption{\textbf{The files and the compute of each sweep.} N9 runs \texttt{run\_n7.py}
    with \texttt{ENGINE\_DTYPE=float32}. I1 and I1b train with \texttt{run\_n7.py} and then
    measure with \texttt{run\_intervention.py}. The hours of a Kaggle sweep are the wall time
    of its T4 session. The hours of a CPU sweep are the sum of the wall times of its runs,
    and some of those runs were in parallel.}
  \label{tab:sweep-files}
  \begin{tabular}{lllll}
    \toprule
    sweep & training script & artifact directory & analysis script & hours \\
    \midrule
    k01 & \texttt{kernels/k01\_scout/run.py} & \texttt{results/k01\_scout} & \texttt{analyze\_scout.py} & $0.6$ \\
    k02 & \texttt{kernels/k02\_grid/run.py} & \texttt{results/k02\_grid} & \texttt{analyze\_k02.py} & $3.7$ \\
    k03 & \texttt{kernels/k03\_grid\_acts/run.py} & \texttt{results/k03\_grid\_acts} & \texttt{analyze\_k03.py} & $3.8$ \\
    k04 & \texttt{kernels/k04\_extended/run.py} & \texttt{results/k04\_extended} & \texttt{analyze\_k04.py} & $3.2$ \\
    k05 & \texttt{kernels/k05\_lowdata/run.py} & \texttt{results/k05\_lowdata} & \texttt{analyze\_k05.py} & $1.3$ \\
    k06 & \texttt{kernels/k06\_horizon/run.py} & \texttt{results/k06\_horizon} & \texttt{analyze\_k06.py} & $4.7$ \\
    k07 & \texttt{kernels/k07\_arma\_seeds/run.py} & \texttt{results/k07\_arma\_seeds} & \texttt{analyze\_k07.py} & $0.5$ \\
    k08 & \texttt{kernels/k08\_init/run.py} & \texttt{results/k08\_init} & \texttt{analyze\_k08.py} & $2.2$ \\
    k09 & \texttt{kernels/k09\_primes\_zdd/run.py} & \texttt{results/k09\_primes\_zdd} & \texttt{analyze\_k09.py} & $1.6$ \\
    N7 & \texttt{run\_n7.py} & \texttt{results/n7\_engine} & \texttt{analyze\_n7.py} & $68.5$ \\
    N9 & \texttt{run\_n7.py} & \texttt{results/n9\_f32} & \texttt{analyze\_n9.py} & $9.6$ \\
    O18 & \texttt{run\_o18\_cpu.py} & \texttt{results/o18\_cpu} & \texttt{analyze\_n9.py} & $1.8$ \\
    O18, f64 & \texttt{run\_o18\_cpu.py -{}-f64} & \texttt{results/o18\_f64} & \texttt{analyze\_n9.py} & $3.2$ \\
    O18-T & \texttt{run\_o18\_transplant.py} & \texttt{results/o18\_transplant} & \texttt{run\_o18\_transplant.py} & $5.3$ \\
    I1 & \texttt{run\_n7.py} & \texttt{results/i1\_internal} & \texttt{run\_intervention.py} & $20.7$ \\
    I1b & \texttt{run\_n7.py} & \texttt{results/i1\_internal} & \texttt{run\_intervention.py} & $24.3$ \\
    Gate 1 & \texttt{run\_gate1.py} & \texttt{results/gate1} & \texttt{test\_gate1.py} & $3.1$ \\
    \bottomrule
  \end{tabular}
\end{table}
% STE end

\subsection{Reproduction procedure}
\label{sec:repro-procedure}

% STE procedural
To reproduce the numbers and the figures of this paper, do these steps:
\begin{enumerate}
  \item Unpack the supplementary archive.
  \item Install Python 3.12.13.
  \item Install the pinned dependencies with \texttt{pip install -r requirements.txt}.
  \item Put the saved artifacts in \texttt{results/} and the run logs in \texttt{logs/}.
  \item From the top directory, run \texttt{bash scripts/reproduce.sh}.
\end{enumerate}
% STE end

% STE descriptive
The script does not train a model. It runs the self-checks and the analysis scripts of
Table~\ref{tab:sweep-files}. The self-checks test Proposition~\ref{prop:torsor} and
Theorem~\ref{thm:stratum} by exhaustion at every modulus up to $200$. The script also runs
\texttt{test\_paper\_numbers.py}, which calculates the numbers in the text, the tables and the
captions again from the artifacts. It fails on a number that it does not calculate, unless
\texttt{paper/number\_ledger.json} records that number with a reason.

\texttt{scripts/render\_all.py} then draws every figure again.
\texttt{scripts/modulus\_table.py} writes Table~\ref{tab:moduli}, and
\texttt{scripts/sweep\_table.py} writes Tables~\ref{tab:sweeps} and~\ref{tab:sweep-files}. The
script needs no GPU. Its longest step is \texttt{analyze\_gate2.py} on the k03 sweep, which
took $8.5$ hours on our CPU.
% STE end

% STE procedural
To train a sweep again, do these steps:
\begin{enumerate}
  \item Run \texttt{scripts/sweep\_table.py -{}-commands}.
  \item Find the sweep. The output gives its git SHA and its exact command.
  \item Run that command from the top directory.
\end{enumerate}
% STE end

% STE descriptive
The self-check of \texttt{sweep\_table.py} makes sure that each command issues the arguments
that the artifacts of its sweep record. It also makes sure that the training code at the
current commit is the same as at that SHA. Where the code is different, the self-check makes
sure that the difference cannot change a result. A Kaggle
sweep runs on a GPU without deterministic settings. Thus its second run is not guaranteed to be
bit-identical to the first (Section~\ref{sec:impl}).
% STE end
